\section{A beta-mixture comparison and the explicit constant}\label{s:beta-mixture}

A block leaves a uniform direction in a random residual subspace. To
bound a high moment of its continuation average, we must control the
next direction after earlier directions in that same subspace have
been exposed. Lemma~\ref{s:directions} describes this law; the next
lemma compares it with two families of tests on which an affine
suffix bound can be used directly.

The first test is the whole ambient sphere. The second fixes one
direction from the span of the earlier draws and averages on a
latitude perpendicular to that direction. If the ambient space has
linear dimension \(D'\), each such latitude has linear dimension
\(D'-1\), regardless of the number of earlier draws. This is the
reason that a block will spend only two test dimensions: one for the
observed projection direction and one for the comparison latitude.
The beta variable in the construction below chooses the latitude
radius so that its law controls the next direction near the preceding
span.

\begin{lemma}[Two-component domination]\label{s:mixture}
Let \(E\) have dimension \(D'\), let \(V\subset E\) have dimension
\(1\le j<n<D'\), and let \(N\) be uniform among the \(n\)-dimensional
subspaces containing \(V\). Conditional on \(N\), let \(v\) be
uniform on \(S(N)\). Set
\[
 \alpha=\frac{n-j}{2},\qquad
 \alpha_*=\frac{D'-j}{2},\qquad
 \beta_0=\frac j2,\qquad
 \beta_1=\frac{D'-n+j-1}{2}.
\]
Let \(\nu_0=\sigma_E\). Let \(\nu_1\) be the law of
\begin{equation}\label{s:mixture-construction}
 \sqrt{1-Z}\,e_0+\sqrt Z\,u_0,\qquad
 Z\sim\Beta(\alpha,\alpha_*-\alpha),\quad
 e_0\sim\sigma_V,\quad
 u_0\mid e_0\sim\sigma_{E\cap e_0^\perp},
\end{equation}
where \(Z\) is independent of the directions. Then, as positive
measures on \(S(E)\),
\begin{equation}\label{s:mixture-domination}
 \Law(v)\le2^{D'}(\nu_0+\nu_1).
\end{equation}
For \(j=0\), the first direction averaged over a uniform \(N\) is
uniform on \(S(E)\). If \(n=D'\), then \(N=E\) and each new
independent direction is uniform even after earlier directions are
conditioned on. These cases require no mixture.
\end{lemma}
\begin{proof}
Put \(Q(y)=\norm{P_{V^\perp}y}^2\). For the actual direction, the
orthogonal decomposition \(N=V\oplus(N\cap V^\perp)\) and
Lemma~\ref{s:coordinates} give
\[
 Q(v)\sim\Beta(\alpha,\beta_0).
\]
For \(\nu_0\), the fraction has law
\(\Beta(\alpha_*,\beta_0)\). Under \(\nu_1\), it equals \(ZX\),
where \(X=\norm{P_{V^\perp}u_0}^2\). If \(j>1\), the decomposition
\[
 E\cap e_0^\perp=V^\perp\oplus(V\cap e_0^\perp)
\]
gives \(X\sim\Beta(\alpha_*,(j-1)/2)\), independently of \(Z\).
The beta-product identity \eqref{s:beta-product} gives
\(Q\sim\Beta(\alpha,\beta_1)\). If \(j=1\), the second summand is
zero and \(X=1\) deterministically. In that case
\(\beta_1=\alpha_*-\alpha\), so \(Q=Z\) has the same stated law.
We do not use a beta law with second shape zero.

The three vector laws are invariant under independent orthogonal
transformations of \(V\) and \(V^\perp\). Their two component lengths
are nonzero almost surely. For completeness, average a bounded Borel
function of a vector over an independent uniform element of
\(O(V)\times O(V^\perp)\). At fixed nonzero component lengths, this
average is integration over the product of the two uniform direction
laws. Invariance leaves the original expectation unchanged, including
after multiplication by any bounded function of \(Q\). Thus each
vector law is obtained from its scalar \(Q\) law with the same
independent uniform component directions. A comparison of the scalar
densities therefore lifts to the full vector measures.

Write \(b=D'+1-n\). The nontrivial case has \(b\ge2\), and
\[
 \alpha_*-\alpha=(b-1)/2>0,\qquad
 \beta_1-\beta_0=(b-2)/2\ge0.
\]
Let \(f_{A,B}\) denote the normalized beta density. The beta integral
decreases when either positive shape parameter increases. For
\(q\ge1/2\),
\[
 \frac{f_{\alpha,\beta_0}(q)}{f_{\alpha_*,\beta_0}(q)}
 =\frac{\mathrm B(\alpha_*,\beta_0)}
        {\mathrm B(\alpha,\beta_0)}
       q^{-(\alpha_*-\alpha)}
 \le2^{\alpha_*-\alpha}\le2^{D'}.
\]
For \(q<1/2\),
\[
 \frac{f_{\alpha,\beta_0}(q)}{f_{\alpha,\beta_1}(q)}
 =\frac{\mathrm B(\alpha,\beta_1)}
        {\mathrm B(\alpha,\beta_0)}
       (1-q)^{-(\beta_1-\beta_0)}
 \le2^{\beta_1-\beta_0}\le2^{D'}.
\]
This includes equality \(\beta_1=\beta_0\) when \(b=2\).
Splitting at \(1/2\) and using the common direction factors proves
\eqref{s:mixture-domination}. For \(j=0\), rotation invariance gives
the uniform marginal; for \(n=D'\), the subspace is the whole \(E\)
and the original draws remain independent uniform directions.
\end{proof}

\begin{proof}[Proof of Theorem~\ref{s:main}]
Fix the memory sequence, and suppose that the conclusion fails along
an unbounded sequence of dimensions. All logarithms in this proof
are to base two. Write
\begin{equation}\label{s:numeric-contradiction}
 \lambda=\log_2(1/\epsilon),\qquad
 c_*=2^{-16},\qquad T<c_*d\lambda.
\end{equation}

\paragraph{The short-run reduction.}
Lemma~\ref{s:short-run}, with \(c=c_*\) and \(\alpha=2/3\), gives
\(\lambda=o(d)\), \(T=o(d^2)\), and
\(\log_2w=o(d^2)\) for the padded width
\[
 w=(T+2)2^M.
\]
These bounds are uniform over the violating accuracies and learner
rules: the explicit estimate \eqref{s:sharp-short-run} bounds
\((d-2)\lambda\) by a function of \(d\) and \(M(d)\) alone.
If \(\lambda\le1024\), then \(T<d/64\), contradicting the
full-data bound \eqref{s:quarter-endpoint} at large \(d\).
We therefore restrict to \(\lambda>1024\).

Set
\begin{equation}\label{s:numeric-parameters}
 k=m=\lfloor d/16\rfloor,\qquad
 p=2\lfloor d/32\rfloor,\qquad
 \ell=p/2,\qquad
 J=\lceil T/k\rceil,\qquad
 G=2^{4d}w^{1/m}.
\end{equation}
The radius exponent \(p\) is even, so the inverse-radius moments
below have the integer exponent \(\ell=p/2\).
Pad early stops to length \(T\) using
Lemma~\ref{s:terminal-count}, and divide the horizon into \(J\)
batches of length at most \(k\). A short last batch retains its actual
positive length. When \(T=0\), there are no batches. At sufficiently
large violating dimensions,
\begin{equation}\label{s:numeric-budget}
 J\le32c_*\lambda+1=o(d),\qquad
 J\le d/8,\qquad m,p\ge d/32.
\end{equation}
The first inequality uses \(k\ge d/32\).

\paragraph{The suffix bound.}
Fix a rule seed for which Lemma~\ref{s:borel} applies, and use its
Borel version. The signal remains uniform and independent of the
pre-generated rows and initialization. Start a continuation from any
specified boundary state, without conditioning on whether that state
is reached. All the bounds below are uniform over this fixed choice
of rules and state.

For a boundary state \(v\) after \(i\) batches, let \(h_{i,v}(s)\)
be its Borel suffix success. We claim that every affine subsphere
\(c+rS(H)\) with \(D=\dim H\ge d-2i\) satisfies
\begin{equation}\label{s:numeric-invariant}
 \int h_{i,v}(c+ru)\,d\sigma_H(u)
 \le G^{J-i}(8\epsilon/r)^p.
\end{equation}
The test dimension in this assertion decreases by two as the
boundary index increases. In the backward step from \(i+1\) to
\(i\), the available bounds therefore include all affine tests of
dimensions \(D-1\) and \(D-2\), whenever the current test has
dimension \(D\ge d-2i\).

At \(i=J\), Lemma~\ref{s:caps} proves the claim. Indeed
\(D\ge d-2J\ge3d/4\) gives \(D-1\ge p\); if
\(8\epsilon/r\ge1\), use total mass, and otherwise compare the
exponents in \eqref{s:affine-terminal}. This also treats \(J=0\).

\paragraph{One block.}
Suppose the claim holds at \(i+1\), and let the next batch have length
\(1\le b\le k\). Lemma~\ref{s:residual} gives
\[
 \rho^2\sim\Beta((D-b)/2,b/2),\qquad
 c'=c+r\sqrt{1-\rho^2}\,e,\qquad r'=r\rho,
\]
\[
 E=H\cap e^\perp,\qquad D'=D-1,\qquad n=D-b.
\]
Conditional on \(\rho,e\), the residual space \(N\) is uniform of
dimension \(n\) in \(E\). By Lemma~\ref{s:state-maximum}, the
left side of \eqref{s:numeric-invariant} is at most
\begin{equation}\label{s:numeric-max}
 \E_{\rho,e}\E_N\max_{v'}F_{v'}(N),\qquad
 F_{v'}(N)=\int h_{i+1,v'}(c'+r'u)\,d\sigma_N(u).
\end{equation}
This step keeps all dependence on the exact batch data until the state
maximum has been taken.

We bound this state maximum through an \(m\)-th moment. Each
conditional factor in that moment will use the same stage-\(i+1\)
affine bound; the earlier replica span changes its distribution but
not the induction stage.

Fix \(\rho,e,v'\) and write
\[
 U_0=G^{J-i-1}(8\epsilon/r')^p.
\]
Expand \(F_{v'}(N)^m\) using \(m\) directions drawn independently
conditional on \(N\). Given the first \(j<m\) directions, their
span \(V\) has dimension \(j\) almost surely, and
Lemma~\ref{s:directions} makes \(N\) uniform among the \(n\)-spaces
containing \(V\). The parameter bounds imply
\[
 D\ge3d/4,\qquad n\ge11d/16,\qquad m\le d/16<n.
\]
For \(j=0\), the next direction is uniform on \(S(E)\), and
the stage \(i+1\) hypothesis bounds its expected success factor by
\(U_0\), since \(D'=D-1\ge d-2(i+1)\). If \(b=1\), then
\(n=D'\) and every next independent direction is uniform in \(E\);
the same bound applies. These cases precede any use of a beta mixture.

For \(j\ge1\) and \(b\ge2\), apply Lemma~\ref{s:mixture}. The
\(\nu_0\) expectation is at most \(U_0\). Under \(\nu_1\),
conditional on \(Z,e_0\), the target is uniform on the affine
subsphere with
\[
 \widetilde H=E\cap e_0^\perp,\qquad
 \widetilde c=c'+r'\sqrt{1-Z}\,e_0,\qquad
 \widetilde r=r'\sqrt Z.
\]
Its linear dimension is \(D-2\ge d-2(i+1)\). Also
\(\widetilde c\perp\widetilde H\), \(\widetilde r>0\), and
\(\norm{\widetilde c}^2+\widetilde r^2=1\), because \(c'\perp E\)
and \(\norm{c'}^2+(r')^2=1\). The induction hypothesis bounds its
mean by \(U_0Z^{-\ell}\).

In the notation of Lemma~\ref{s:mixture},
\[
 n-j\ge10d/16,\qquad
 \alpha-\ell\ge9d/32\ge d/4,\qquad
 \alpha_*\le d/2.
\]
The other shape \(\alpha_*-\alpha=(b-1)/2\) is positive.
The integer \(\ell=\lfloor d/32\rfloor\) is strictly below the
first shape, and every ratio below is at most two:
\begin{equation}\label{s:mixture-inverse}
 \E Z^{-\ell}
 =\prod_{t=1}^{\ell}\frac{\alpha_*-t}{\alpha-t}
 \le2^\ell\le2^d.
\end{equation}
Since \(D'\le d\), the mixture domination bounds the conditional
expected next factor by
\[
 2^dU_0(1+\E Z^{-\ell})\le2^{3d}U_0.
\]
For each preceding replica span, this uses a fresh test in the same
dimension \(D-2\). Only the single direction \(e_0\) is removed from
\(E\); no dimension decrease is carried from one replica to the next.

Integrating the nonnegative product in reverse conditional order gives
\[
 \E_NF_{v'}(N)^m\le(2^{3d}U_0)^m.
\]
Jensen's inequality and the sum of the at most \(w\) state moments
then give
\[
 \E_N\max_{v'}F_{v'}(N)\le w^{1/m}2^{3d}U_0.
\]
For the residual fraction \(\rho^2\), the first shape satisfies
\(n/2-\ell\ge5d/16\ge d/4\), its total shape is \(D/2\le d/2\),
and its other shape \(b/2\) is positive even for a one-row batch.
The same integer beta identity yields
\begin{equation}\label{s:residual-inverse-numeric}
 \E\rho^{-p}
 =\prod_{t=1}^{\ell}\frac{D/2-t}{n/2-t}
 \le2^\ell\le2^d.
\end{equation}
Averaging \(r'=r\rho\) in \eqref{s:numeric-max} supplies the factor
\(w^{1/m}2^{4d}=G\), completing the backward induction.

\paragraph{The numerical conclusion.}
At \(i=0\), take the full sphere \(H=\R^d,c=0,r=1\).
The bound is uniform in the fixed shared seed and initial state;
joint independence permits averaging over initialization, and then
Lemma~\ref{s:borel} permits averaging the original measurable
conditional success probabilities over the seed. By
\(\log_2w=o(d^2)\) and \(m\ge d/32\), we have
\(\log_2G\le5d\) eventually. The base-two logarithm of the success
upper bound is at most
\begin{align}
 5d(32c_*\lambda+1)+p(3-\lambda)
 &\le5d+(160c_*-1/64)d\lambda\notag\\
 &\le5d-d\lambda/128\le-3d.\label{s:numeric-arithmetic}
\end{align}
Here \(3-\lambda\le-\lambda/2\), \(p\ge d/32\),
\(160c_*=5/2048\le1/128\), and \(\lambda\ge1024\).
The resulting success is less than \(2/3\), a contradiction.

Both precision branches exclude an unbounded sequence of violations.
All eventual conditions were bounded using
\eqref{s:sharp-short-run} and the fixed memory sequence; the
full-data cap threshold is absolute. This proves the uniform
dimension threshold and \eqref{s:main-bound}.
\end{proof}
